% ============================================================
% Multi-Domain Severity Scales — Empirical Calibration
% For EMNLP 2026 main text or appendix
% ============================================================

\section{Multi-Domain Severity Scales}
\label{sec:multi-domain-scales}

We define domain-specific cost vectors $\mathbf{c} = (c_1, c_2, c_3, c_4)$
calibrated against publicly available regulatory, litigation, and
market data. Each domain uses the same four severity labels
(\textit{negligible}, \textit{minor}, \textit{major}, \textit{critical})
but with domain-appropriate dollar amounts.

% ============================================================
\subsection{Finance (\$100, \$1k, \$10k, \$100k)}

See Section~\ref{sec:severity-scale} for full calibration.

\textbf{Key sources:} SEC SAB~99 materiality framework; FINRA median
fine \$125k (2024); SOX~\S906 individual penalty \$10k; ``Big~R''
restatement average net income impact --\$16M; IBM/Ponemon cost per
record \$173.

% ============================================================
\subsection{Medical (\$500, \$5k, \$50k, \$500k)}

Medical errors carry higher unit costs than financial errors due to
direct patient safety implications and the regulatory environment
(FDA, malpractice liability).

\begin{table}[h]
\centering
\small
\begin{tabular}{lrl}
\toprule
Level & $c_k$ & Calibration \\
\midrule
Negligible & \$500 & BMI/BSA miscalculation; unit conversion \\
Minor & \$5\,000 & Framingham risk score error; non-critical lab \\
Major & \$50\,000 & Surgical risk miscalculation; delayed diagnosis \\
Critical & \$500\,000 & Dosage error; sepsis triage failure \\
\bottomrule
\end{tabular}
\caption{Medical severity scale.}
\label{tab:medical-scale}
\end{table}

\paragraph{Empirical evidence.}
\begin{itemize}
\item \textbf{AHRQ/HCUP 2023:} Average cost of a preventable adverse
    drug event (ADE): \$5\,857 per hospital admission
    \citep{ahrq2023ade}. Calibrates \textit{minor}.
\item \textbf{Johns Hopkins/BMJ 2016:} Medical errors are the third
    leading cause of death in the U.S., with estimated annual cost of
    \$17--29B \citep{makary2016}. Calibrates systemic exposure.
\item \textbf{Malpractice settlements:} Median jury verdict in medical
    malpractice: \$1.1M; median settlement: \$425k
    \citep{diederich2023}. Calibrates \textit{critical}.
\item \textbf{FDA medication error reports:} 100\,000+ reports/year;
    7\,000--9\,000 deaths attributed to medication errors annually
    \citep{fda2023mederrors}. Per-death economic cost estimate:
    \$1--10M (value of a statistical life, EPA/DOT).
\item \textbf{Diagnostic error cost:} Estimated \$100B/year in the U.S.
    from diagnostic errors; average per-patient cost of a missed
    diagnosis: \$36\,000--\$85\,000 \citep{newman2023diagnostic}.
    Calibrates \textit{major}.
\item \textbf{Sepsis misdiagnosis:} Delayed sepsis treatment increases
    mortality by 7.6\% per hour of delay; hospital cost per sepsis
    case: \$32\,000--\$48\,000; mortality-adjusted cost exceeds \$500k
    \citep{seymour2017sepsis}. Calibrates \textit{critical}.
\end{itemize}

\paragraph{Annotation rule for MedCalc-Bench.}
The dataset categorizes medical calculators into clinical domains
(dosage, triage, risk assessment, body metrics). The mapping is:
\begin{itemize}
\item Dosage/Sepsis/Triage calculators $\to$ \textit{critical}
\item Surgical risk / APACHE / ICU scores $\to$ \textit{major}
\item Framingham / cardiovascular risk scores $\to$ \textit{minor}
\item BMI / BSA / unit conversions $\to$ \textit{negligible}
\end{itemize}

% ============================================================
\subsection{Legal (\$200, \$2k, \$20k, \$200k)}

Legal errors in contract analysis expose organizations to missed
obligations, unintended liability, and unfavorable terms.

\begin{table}[h]
\centering
\small
\begin{tabular}{lrl}
\toprule
Level & $c_k$ & Calibration \\
\midrule
Negligible & \$200 & Definitions, boilerplate identification \\
Minor & \$2\,000 & Effective date, term, renewal errors \\
Major & \$20\,000 & Non-compete, IP assignment, exclusivity \\
Critical & \$200\,000 & Indemnification, liability cap, termination \\
\bottomrule
\end{tabular}
\caption{Legal severity scale.}
\label{tab:legal-scale}
\end{table}

\paragraph{Empirical evidence.}
\begin{itemize}
\item \textbf{IACCM/WCC 2022:} Average cost of poor contract management:
    9.2\% of annual revenue \citep{iaccm2022}. For a \$100M company,
    this is \$9.2M/year across all contracts.
\item \textbf{Legal billing rates:} Average partner billing rate
    \$600--1\,200/hour (Am Law 200). A missed clause requiring
    2--20 hours of remediation: \$1\,200--\$24\,000.
    Calibrates \textit{minor} to \textit{major}.
\item \textbf{Indemnification disputes:} Median cost of contractual
    indemnification litigation: \$200k--\$500k including legal fees
    \citep{aba2021indemnification}. Calibrates \textit{critical}.
\item \textbf{Non-compete violations:} Average damages awarded in
    non-compete cases: \$18\,000--\$75\,000 for injunctive relief;
    \$100k--\$500k for lost profits claims. Calibrates
    \textit{major} to \textit{critical}.
\item \textbf{M\&A due diligence errors:} A single missed material
    contract clause can affect deal valuation by 1--5\% of
    transaction value. For a \$50M deal: \$500k--\$2.5M.
\end{itemize}

\paragraph{Annotation rule for CUAD.}
CUAD annotates 41 contract clause types. The mapping is based on the
legal risk category:
\begin{itemize}
\item Indemnification, limitation of liability, uncapped liability,
    termination for convenience $\to$ \textit{critical}
\item Non-compete, IP ownership, exclusivity, most favored nation,
    change of control $\to$ \textit{major}
\item Effective date, expiration date, renewal term, notice period,
    governing law $\to$ \textit{minor}
\item Document name, parties, definitions, recitals $\to$
    \textit{negligible}
\end{itemize}

% ============================================================
\subsection{Code/Security (\$1k, \$10k, \$100k, \$1M)}

Vulnerability-related errors in code carry costs that scale with
the CVSS severity score and the attack surface affected.

\begin{table}[h]
\centering
\small
\begin{tabular}{lrl}
\toprule
Level & $c_k$ & Calibration \\
\midrule
Negligible & \$1\,000 & Low CVSS (0.1--3.9); informational finding \\
Minor & \$10\,000 & Medium CVSS (4.0--6.9); limited exploit \\
Major & \$100\,000 & High CVSS (7.0--8.9); network-exploitable \\
Critical & \$1\,000\,000 & Critical CVSS (9.0--10.0); RCE/zero-day \\
\bottomrule
\end{tabular}
\caption{Code/Security severity scale.}
\label{tab:code-security-scale}
\end{table}

\paragraph{Empirical evidence.}
\begin{itemize}
\item \textbf{IBM/Ponemon 2024:} Average cost of a data breach:
    \$4.88M globally; \$9.48M in the U.S. \citep{ponemon2024}.
    Average cost per stolen record: \$173. A critical vulnerability
    leading to a breach of 5\,000--50\,000 records: \$865k--\$8.65M.
    Calibrates \textit{critical}.
\item \textbf{Bugcrowd 2024:} Median bug bounty payout for critical
    vulnerabilities: \$8\,000--\$30\,000; for high: \$2\,000--\$5\,000
    \citep{bugcrowd2024}. Note: bounties reflect the \emph{research}
    cost, not the \emph{impact} cost, providing a lower bound.
\item \textbf{NIST NVD statistics:} 28\,902 CVEs published in 2023;
    17.8\% rated Critical (CVSS $\geq$ 9.0); 42.3\% rated High
    \citep{nist2024nvd}.
\item \textbf{Equifax breach (2017):} Unpatched Apache Struts
    vulnerability (CVSS~10.0); total cost: \$1.4B including \$700M
    settlement \citep{ftc2019equifax}. Single CVE, \$1.4B impact.
\item \textbf{Log4Shell (2021):} CVSS~10.0; estimated global
    remediation cost: \$10B+ \citep{cybereason2022log4shell}.
\item \textbf{SolarWinds (2020):} Estimated cleanup cost across
    affected organizations: \$100M+ per large enterprise
    \citep{solarwinds2021}.
\item \textbf{Veracode State of Software Security 2024:} 74\% of
    applications have at least one security flaw; 19.5\% have
    high-severity flaws \citep{veracode2024}.
\end{itemize}

\paragraph{Annotation rule for CVE-Bench.}
CVE-Bench provides native CVSS scores. The mapping is direct:
\begin{itemize}
\item CVSS $\geq$ 9.0 $\to$ \textit{critical}
\item CVSS 7.0--8.9 $\to$ \textit{major}
\item CVSS 4.0--6.9 $\to$ \textit{minor}
\item CVSS $<$ 4.0 $\to$ \textit{negligible}
\end{itemize}

% ============================================================
\subsection{Content Moderation (\$50, \$500, \$5k, \$50k)}

Content moderation errors expose platforms to regulatory fines
(DSA, Online Safety Act), advertiser withdrawal, and user harm.

\begin{table}[h]
\centering
\small
\begin{tabular}{lrl}
\toprule
Level & $c_k$ & Calibration \\
\midrule
Negligible & \$50 & Borderline content; minor policy edge case \\
Minor & \$500 & Hate speech / harassment missed \\
Major & \$5\,000 & CSAM-adjacent; terrorist content missed \\
Critical & \$50\,000 & CSAM; imminent violence threat missed \\
\bottomrule
\end{tabular}
\caption{Content moderation severity scale.}
\label{tab:moderation-scale}
\end{table}

\paragraph{Empirical evidence.}
\begin{itemize}
\item \textbf{EU Digital Services Act (2024):} Fines up to 6\% of
    global annual turnover for systemic failures in content moderation.
    For a platform with \$10B revenue: up to \$600M.
\item \textbf{UK Online Safety Act (2023):} Fines up to 10\% of
    qualifying worldwide revenue, or \pounds18M (whichever is higher).
\item \textbf{Australia eSafety Commissioner:} AUD~555k per day per
    contravention for failure to remove class~1 material (CSAM,
    terrorism).
\item \textbf{Meta Oversight Board:} Reports that incorrect content
    removal affects an estimated 1 in 10\,000 pieces of content;
    incorrect non-removal of violating content estimated at 3--5\%
    of reports \citep{meta2023transparency}.
\item \textbf{Advertiser boycotts:} The 2020 ``Stop Hate for Profit''
    campaign cost Facebook an estimated \$7.2B in market value in
    one week \citep{reuters2020stophate}.
\item \textbf{Trust \& Safety teams:} Average cost per manual content
    review: \$1--5 per item (including human reviewer time, QA,
    appeals). At scale (millions of reviews/day), this represents
    the \textit{negligible} operational baseline.
\end{itemize}

\paragraph{Annotation rule for Jigsaw Severity.}
Jigsaw provides pairwise severity comparisons. We convert the
ranking-derived severity score to quartiles:
\begin{itemize}
\item Top 5\% severity $\to$ \textit{critical}
\item 75th--95th percentile $\to$ \textit{major}
\item 25th--75th percentile $\to$ \textit{minor}
\item Bottom 25\% $\to$ \textit{negligible}
\end{itemize}

% ============================================================
\subsection{Summary of Cost Vectors}

\begin{table}[h]
\centering
\small
\begin{tabular}{lrrrr}
\toprule
Domain & $c_1$ & $c_2$ & $c_3$ & $c_4$ \\
\midrule
Finance         & \$100   & \$1\,000   & \$10\,000   & \$100\,000 \\
Medical         & \$500   & \$5\,000   & \$50\,000   & \$500\,000 \\
Legal           & \$200   & \$2\,000   & \$20\,000   & \$200\,000 \\
Code/Security   & \$1\,000 & \$10\,000 & \$100\,000  & \$1\,000\,000 \\
Moderation      & \$50    & \$500      & \$5\,000    & \$50\,000 \\
\bottomrule
\end{tabular}
\caption{Cost vectors by domain. Each vector spans three orders of
magnitude, reflecting the empirical observation that financial exposure
from errors follows a heavy-tailed distribution within each domain.}
\label{tab:cost-vectors}
\end{table}

\paragraph{Cross-domain ordering.}
The relative ordering $c_k^{\text{moderation}} < c_k^{\text{finance}}
< c_k^{\text{legal}} < c_k^{\text{medical}} < c_k^{\text{security}}$
reflects the regulatory and liability gradient: content moderation
errors rarely trigger direct financial liability at the per-instance
level, while a single critical vulnerability (RCE/zero-day) can lead
to breaches costing hundreds of millions.
